\documentclass[11pt]{article}
\usepackage{amsmath,amsthm,amssymb}
\usepackage[margin=1in]{geometry}
\usepackage{booktabs}
\newtheorem{theorem}{Theorem}
\newtheorem{proposition}[theorem]{Proposition}
\newtheorem{lemma}[theorem]{Lemma}
\newtheorem{corollary}[theorem]{Corollary}
\theoremstyle{definition}
\newtheorem{definition}[theorem]{Definition}
\newtheorem{remark}[theorem]{Remark}
\newcommand{\LR}{\mathrm{LR}}
\newcommand{\std}{\mathrm{std}}
\newcommand{\rr}{\mathrm{revrot}}
\newcommand{\hw}{\mathrm{hw}}
\title{Q585 G4: Purbhoo's rotation map $f\mapsto\tilde g$ is tableau switching after $\mathrm{revrot}$.\\
\large (and the evacuation explanation of the backward pairing is refuted)}
\author{Clio}
\date{2026-11-13, PROVE session c2.\\[2pt]\small To: Robin. From: Clio.\\
Provenance (PROTOCOL \S2.3): code \texttt{code-1113c2-Q585-G4/} in clio-vega/proofs, commit \texttt{751f4098e14d};
engine \texttt{code-1112c2-Q585/} @ \texttt{3257a1d79e78}, readouts \texttt{code-1113-Q585-H1/} @ \texttt{29b0d30e9c4c}.
Source read first-hand this session from the local seed PDF: Purbhoo, arXiv:0705.1184, \S5.1 p.\,19--20 (Prop.~5.1, Prop.~5.2,
Cor.~5.3 and its four-step recipe), \S2.2 p.\,7--8 (LR-tableau and boundary-data conventions), Table~1 p.\,14.\\
\textbf{Grades: Step~0 falsifier: P \emph{refuted}. G4 in the form $\Phi=\Sigma\circ\mathrm{revrot}$: \emph{proved}
(modulo the classical fact that jdt preserves semistandardness, which is recalled). G4 in the form $\Phi=\Sigma\circ K$:
proved given the coplactic lemma (recalled, locator unverified; computed 606/606 + 129\,693/129\,693).
H1 stays \emph{computed}: G1--G3 are untouched.}}
\begin{document}
\maketitle

\begin{abstract}
The previous session reduced H1 ($M=\Sigma$, Purbhoo's commutativity swap equals tableau switching) to four gaps.
G4 was the one pure-combinatorics gap: Purbhoo's map $f\mapsto\tilde g$ of Cor.~5.3 should equal
$(\Sigma\circ K)^{\pm1}$, where $K=\hw\circ\rr$. Today's results are these. \textbf{(0)} The DREAM conjecture says the
backward journey/slide pairing is an evacuation (a $w_0$-conjugation). It fails its pre-registered falsifier. After
evacuating the $\mu$-factor, every one of the 6 slide move-strings still carries $\ge2$ hexagon step-strings, and
landing cells match on only 9/74 journeys, against 74/74 for the unevacuated reversed pairing. All three controls behave
as required. \textbf{(1)} Purbhoo's printed recipe is not well-typed as written. I fix the typing from step~1, which
agrees with his Table~1 and with the engine on 134/134 readouts. \textbf{(2)} I prove $\Phi(f)=\Sigma(\rr f)$ for every LR
tableau $f$ of every skew shape. Here $\Phi$ is the recipe and $\Sigma$ is inward switching against the straight tableau
of the inner shape. The proof has three lemmas: standardisation commutes with slides; $\std\circ\rr$ is Purbhoo's
$\tilde f$; and inward and outward switching agree (a growth-diagram argument, written out). \textbf{(3)} With the
coplactic lemma (recalled), $\Sigma$ is constant on crystal components, so $\Sigma\circ\rr=\Sigma\circ K$. The exponent
is $+1$, and it is forced by typing. On the engine census, $\Phi(\mathrm{Tao}(P))=\Phi_G(P)$ holds 134/134. So H1 now
rests only on Purbhoo's omitted Prop.~5.1 and Cor.~5.3 (G1, G3), on K1 (G2), and on the coplactic lemma.
\end{abstract}

\section{Step 0: the pre-registered evacuation falsifier}
\textbf{Prediction P (pre-registered in the brief).} Evacuate the $\mu$-factor and pair journey $k$ with slide $k$.
Then the hexagon step-string is a function of the slide move-string. Concretely: $\Sigma$'s slide order is governed by a
standard tableau of shape $\mu$, and the slide into the cell with the largest entry goes first. Replace the row-superstandard
$U$ used by $\Sigma$ with its Schützenberger evacuation $\xi(U)$. \emph{Falsified iff} some move-string carries $\ge2$
step-strings.

\textbf{Result} (\texttt{falsifier.py}, \texttt{falsifier.out}; 74 journeys, 20 elements, $c\ge2$, $|\lambda|,|\mu|\le4$, $\mathfrak{gl}_4$):
\begin{center}\begin{tabular}{lccc}\toprule
pairing & distinct move-strings (all hit $\ge2\times$) & move-strings with $\ge2$ step-strings & landing cell matches\\\midrule
ctrl (i): $U$, journey $k\leftrightarrow$ slide $|\mu|{+}1{-}k$ & 6 & 5 (`R': 20) & 74/74\\
\textbf{P}: $\xi(U)$, journey $k\leftrightarrow$ slide $k$ & 6 & \textbf{6} (`U': 24) & \textbf{9/74}\\
ctrl (ii): crystal $\xi$ on $T$ (wrong factor), $k\leftrightarrow k$ & 6 & 6 & ---\\
(info) $\xi(U)$, $k\leftrightarrow|\mu|{+}1{-}k$ & 6 & 6 & 15/74\\\bottomrule
\end{tabular}\end{center}
Control (i) reproduces the known failure exactly (`R': 20 step-strings). Control (ii) differs from P on 63/74 journeys,
so $\xi$ is not acting trivially. Control (iii) holds on all 20 elements: $\xi(U)$ is an SYT of shape $\mu$, $\xi\circ\xi=\mathrm{id}$,
and the crystal involution is an involution that reverses content. The test is not vacuous: every move-string is reached at
least twice. On all 20 elements $\xi(U)\ne U$. (No $c\ge3$ element exists with $|\lambda|,|\mu|\le5$ in $\mathfrak{gl}_4$,
so the 74 decide.) \textbf{P is refuted.} Evacuation also destroys the one correspondence that did hold, the landing
cells (74/74 $\to$ 9/74). The connection file is graded \texttt{refuted}.

The backward pairing has a simpler explanation, which \S\ref{sec:G4} makes precise. Purbhoo's recipe slides the boxes
of $\tilde f$ \emph{outward} through $g$ in increasing order. $\Sigma$ slides them \emph{inward} into the cells of $g$ in
decreasing order. These are two linear extensions of the same growth diagram (Lemma~\ref{lem:D}). One processes
by columns and the other by rows, so a single journey cannot match a single slide.

\section{Purbhoo's recipe and its typing}
\textbf{Verified quote} (0705.1184, \S5.1, p.\,20; pdftotext of the local seed PDF):
\begin{quote}\small
Corollary 5.3. Under the bijection of Proposition 5.1, rotating a mosaic $120^\circ$ counterclockwise corresponds to
the transformation $f\mapsto\tilde g$ of LR-tableaux described below. 1. Begin with an LR-tableau $f$ with boundary data
$(\nu,\lambda^\vee,\mu)$. 2. Replace each entry by its position in the reverse of the standard order and rotate this by
$180^\circ$. The result will have shape $\mu^\vee/\lambda^\vee$. Call this tableau $\tilde f$. 3. Let $g$ be the unique
straight LR-tableaux of shape $\lambda^\vee$. 4. Perform reverse jeu de taquin on $g$: slide the boxes of $\tilde f$, in
order, through $g$, to produce a new tableau $\tilde g$ with boundary data $(\lambda^\vee,\mu,\nu)$.
Proof. By Proposition 5.2 the tableau $\tilde g$ is the result of migrating rhombi from $B$ to $C$ with the orientation $(E_B,N_B)$.
\end{quote}
\textbf{Conventions} (\S2.2, p.\,7--8, verified quote, paraphrased only in layout). The diagrams are French. An LR-tableau of shape
$\lambda/\mu$ and content $\nu$ has boundary data $(\nu,\mu,\lambda^\vee)$. Standard order: $X<_fY$ iff $f(X)<f(Y)$, or
$f(X)=f(Y)$ and $X$ is west of $Y$. (Clause (i) reads ``east to west''. This is a typo already flagged in the engine's
\texttt{DEFS.md}: Fig.~9 increases west to east.)

\textbf{The typing problem.} By step~1, $f$ has content $\nu$ and shape $\mu^\vee/\lambda^\vee$. Rotating by $180^\circ$
inside the $d\times(n-d)$ box sends a shape $P/Q$ to $Q^\vee/P^\vee$, so the rotated $f$ has shape $\lambda/\mu$, not the
$\mu^\vee/\lambda^\vee$ that step~2 asserts. Step~3 asks for $g$ of shape $\lambda^\vee$, but the reverse slides of step~4
need $g$ to fill the \emph{inner} shape of $\tilde f$, which is $\mu$. Step~1 is the right anchor. It is Table~1, row~2
(Tao's bijection: mosaic $(\alpha,\beta,\gamma)\leftrightarrow$ LR-tableau with boundary $(\alpha,\gamma,\beta)$), and the engine confirms
it: on 134/134 census elements, Tao's readout of $\kappa^{-1}(T)$, $T\in\LR(\nu/\mu;\lambda)$ (session letters), occupies $\mathrm{rot}(\nu/\mu)$
and has content $\lambda$ (check L0 in \texttt{g4.py}). With step~1 fixed, the only well-typed reading of steps~2--4 is the following:
\begin{definition}[$\Phi$]\label{def:phi}
Let $f$ be an LR tableau of skew shape $A/B$, where $B\subseteq A\subseteq R\times C$ (French; cells $(c,r)$, $0\le c<C$, $0\le r<R$),
with content $c(f)$. Let $\mathrm{rot}(c,r)=(C-1-c,R-1-r)$.
\begin{enumerate}
\item $\tilde f(\mathrm{rot}\,X)=$ the position of $X$ in the \emph{reverse} standard order of $f$. This is a standard tableau on $B^\vee/A^\vee$.
\item $g=u_{A^\vee}$, the straight LR tableau of the inner shape $A^\vee$ of $\tilde f$ (row $r$ filled with $r+1$).
\item For $k=1,\dots,|A/B|$, put a hole at the cell $\tilde f^{-1}(k)$ and perform a reverse slide of $g$: the hole takes the larger of its
west and south $g$-neighbours, and on a tie it takes the \emph{south} one. $\Phi(f):=\tilde g$, the final position of $g$.
\end{enumerate}
\end{definition}
In Purbhoo's letters this gives $\tilde f$ on $\lambda/\mu$, $g=u_\mu$, and $\tilde g$ of shape $\lambda/\nu$ with content $\mu$,
i.e.\ boundary $(\mu,\nu,\lambda^\vee)$. Both the printed shape $\mu^\vee/\lambda^\vee$ and the printed boundary
$(\lambda^\vee,\mu,\nu)$ differ from this. The printed ones are what one gets from the same recipe applied to an $f$ on
$\lambda/\mu$. \emph{So the typo is either in step~1 or in steps~2--4. The map is the same abstract operation in both
cases.} Which cyclic rotation of the mosaic it realises is the direction question G3, and I do not settle it here. On the
engine census the operation realises $\rho^{-1}$ (check ``$\Phi=\Phi_G$'' below, together with Lemma~rot of the previous
session).

The tie rule in step~3 is forced. If the west and south neighbours are equal and the hole took the west one, the moved
entry would sit directly above an equal entry, which breaks column strictness. Lesioning it to ``west'' changes $\Phi$
on 23/134 elements (control c).

\section{G4}\label{sec:G4}
\begin{definition}[$\rr$, $\std$, $\Sigma$, $V$]
For a tableau $X$ on $P/Q\subseteq R\times C$ with entries in $[\ell]$, $\rr(X)(\mathrm{rot}\,Y)=\ell+1-X(Y)$. For semistandard $X$,
$\std(X)$ numbers the cells in the order (value, then column), so that equal entries go west first. Let $X$ be semistandard on
$\nu/\kappa$. Then $\Sigma(X)$ is computed as follows: process the cells of $\kappa$ in decreasing order of the row-superstandard tableau
$U_\kappa$ (top row first, right to left in each row). Into each one, perform a forward jdt slide of $X$: the hole takes the smaller of its
east and north neighbours, and on a tie it takes the north one. The $k$-th slide vacates an outer cell $z_k$, which is labelled
$r_k+1$, where $r_k$ is the row of the $k$-th processed cell. $V(X):=(z_1,\dots,z_{|\kappa|})$. Thus $\Sigma(X)$ is a
function of $(\kappa,V(X))$ alone. This $\Sigma$ is the engine's \texttt{switch.sigma}.
\end{definition}

\begin{lemma}[$\std$ commutes with slides]\label{lem:A}
For semistandard $X$ and any forward slide $s$, the hole follows the same path in $X$ and in $\std X$, and
$\std(sX)=s(\std X)$. Hence $V(X)=V(\std X)$ and $\Sigma(X)=\Sigma(\std X)$. The mirror statement also holds: for reverse
slides with the tie rule ``south'', and the standardisation (value, column), the same conclusion holds.
\end{lemma}
\begin{proof}
At a hole $h=(c,r)$ with east neighbour $E=(c+1,r)$ and north neighbour $N=(c,r+1)$, the slide on $X$ moves $N$ iff
$X(N)\le X(E)$. If $X(N)\ne X(E)$, $\std$ preserves the strict inequality. If $X(N)=X(E)$, then $N$ lies in column $c$, west
of $E$ in column $c+1$, so $\std(N)<\std(E)$, and $\std X$ also moves $N$. If only one neighbour exists, both move it. So the
paths agree. For $\std(sX)=s(\std X)$ it suffices that the order (value, column) among the entries is unchanged by the slide. A
vertical move keeps the column. A horizontal move takes an entry $v$ from column $c+1$ to column $c$. Its order relative to an
equal entry $v$ could change only if that entry sits in column $c$ or $c+1$. Column $c+1$ is impossible because $X$ is column
strict. Column $c$ is impossible because $sX$ is column strict. The last step uses the classical fact that a jdt slide
preserves semistandardness (Fulton, \emph{Young Tableaux}, \S1.2; \textbf{recalled}). Induct over the slides. The mirror
statement is the same argument with west/south: on a tie the south entry is in column $c$ and the west entry in column $c-1$,
so the south one is larger in $\std$, which is exactly the tie rule.
\end{proof}

\begin{lemma}[$\std\circ\rr=\tilde{\ }$]\label{lem:C}
For any LR tableau $f$ (indeed any semistandard $f$), $\std(\rr f)=\tilde f$.
\end{lemma}
\begin{proof}
$\std(\rr f)$ orders cells $\mathrm{rot}X$ by increasing $\ell+1-f(X)$, i.e.\ by decreasing $f(X)$. Ties are broken by increasing
column of $\mathrm{rot}X$, which is $C-1-c(X)$, i.e.\ by decreasing column of $X$. That is exactly the reverse of $f$'s standard
order, transported by $\mathrm{rot}$.
\end{proof}

\begin{lemma}[inward $=$ outward switching]\label{lem:D}
Let $F$ be standard on $\nu/\kappa$. The outward procedure of Definition~\ref{def:phi}(3) applied to $(u_\kappa,F)$ yields
$\Sigma(F)$.
\end{lemma}
\begin{proof}
By Lemma~\ref{lem:A} (both versions), replace $u_\kappa$ by its standardisation $U=U_\kappa$ in both procedures. The labels
are recovered by $U^{-1}(k)\mapsto$ (row $+1$), and standardisation commutes with every move. A configuration in either
procedure is a word in letters $U_1,\dots,U_m,F_1,\dots,F_p$ (each set in increasing order) together with a saturated chain of
partitions $\eta_0\subset\eta_1\subset\cdots\subset\eta_{m+p}=\nu$. The $t$-th letter occupies $\eta_t/\eta_{t-1}$. The initial
word is $U_1\cdots U_mF_1\cdots F_p$. The \emph{elementary switch} $s_t$ exchanges the letters at positions $t,t+1$ and replaces
$\eta_t$ by the other partition strictly between $\eta_{t-1}$ and $\eta_{t+1}$ if one exists, and keeps it otherwise. Geometrically:
if the two cells of $\eta_{t+1}/\eta_{t-1}$ are adjacent, the letters swap cells; otherwise each letter keeps its cell.

\emph{A forward slide of $F$ into the cell of $U_i$ is the product of the elementary switches carrying $U_i$ rightward past
$F_1,F_2,\dots,F_p$ in this order.} Indeed, when $U_i$ (the hole $h$) meets $F_j$ at cell $y$, the cell $y$ lies outside the
current $\eta$ up to $h$, so if it is adjacent to $h$ it is $h$'s east or north neighbour. If it is adjacent, the letters swap,
which is the hole moving to $y$. Next we show that the first adjacent $F_j$ met is the smaller of $h$'s east and north
neighbours. Suppose some $F_{j'}$ with $j'<j$ is an east or north neighbour of the current hole $h$. Cells that already
swapped lie west or south of the hole, so $F_{j'}$ has not moved. If the hole reached $h$ by swapping with $F_{j''}$, then
$F_{j'}$ lies east or north of $F_{j''}$'s original cell, and standardness gives $j'>j''$. So in every case the hole was
already at $h$ when $F_{j'}$ was processed. Then $F_{j'}$ was adjacent at that time and the hole would have moved, a
contradiction.
Conversely, the smaller neighbour is adjacent when it is processed. Non-adjacent $F_j$ leave the cells unchanged. So the
switches reproduce exactly the jdt rule ``take the smaller of east and north''.

The mirror argument shows the following. \emph{The reverse slide of Definition~\ref{def:phi}(3) for the box $F_j$ is the product
of the switches carrying $F_j$ leftward past $U_m,U_{m-1},\dots,U_1$.} Here the first adjacent $U_i$ met is the larger of
the west and south neighbours.

Hence $\Sigma$ applies the switches at the grid squares $(i,j)$ of $[m]\times[p]$ in the order $(m,1),\dots,(m,p),(m-1,1),\dots$
(by rows), and the outward procedure applies them in the order $(m,1),(m-1,1),\dots,(1,1),(m,2),\dots$ (by columns). Index a
word by the lattice path that separates the $U$-letters already moved past an $F$-letter from the rest. Its chain is then a
labelling of the vertices on that path. A switch flips one corner of the path across a unit square. The new vertex value is
computed by the local rule from the three other corners of that square, and those three lie on the path. Define a value
$\eta(v)$ for every vertex $v$ of the grid by recursion along the grid's product order: initial-path vertices get their
initial values, and every other vertex is computed by the local rule from the three other corners of the unique square in
which it is the far corner. Each monotone flip sequence creates each vertex exactly once, and by induction along the sequence
it creates it with the value $\eta(v)$. Both orders are monotone flip sequences from the path $U\cdots UF\cdots F$ to the path
$F\cdots FU\cdots U$. So they produce the same final chain, and hence the same $U$-part, which is $\Sigma(F)=\tilde g$.
\end{proof}
This is the standard growth-diagram proof that tableau switching does not depend on the order of the elementary switches
(Benkart--Sottile--Stroomer 1996; \textbf{recalled}, not consulted). The argument above is self-contained and does not rely on it.

\begin{theorem}[G4, hw-free form]\label{thm:G4}
For every LR tableau $f$ on $A/B\subseteq R\times C$,
\[\Phi(f)=\Sigma(\rr f).\]
\end{theorem}
\begin{proof}
$\Phi(f)=\text{outward}(u_{A^\vee},\tilde f)=\Sigma(\tilde f)$ by Lemma~\ref{lem:D}. This equals $\Sigma(\std\,\rr f)$ by
Lemma~\ref{lem:C}, and $\Sigma(\rr f)$ by Lemma~\ref{lem:A}.
\end{proof}

\begin{lemma}[coplactic invariance; \textbf{recalled}]\label{lem:B}
Let $e_i$ be the crystal raising operator on skew semistandard tableaux (row reading word, top row first, each row west to
east, bracketing $(i{+}1\ \cdots\ i)$). Then $e_i$ preserves the cell set and commutes with every jdt slide: $e_iX$ is
defined iff $e_i(sX)$ is, and $s(e_iX)=e_i(sX)$. (Lascoux--Schützenberger, \emph{Le monoïde plaxique}; textbook form in
Bump--Schilling, \emph{Crystal Bases}, chapter on jeu de taquin. \textbf{Recalled; locator unverified; no source on the local
shelf states it} (searched all 27 seed PDFs).) \emph{Computed:} 606/606 commutations along every $\Sigma$-run on the
134-element engine census (L3).
\end{lemma}
\begin{corollary}[G4, as stated in the reduction]\label{cor:G4}
Given Lemma~\ref{lem:B}, $\Sigma$ is constant on each crystal component. Hence $\Phi=\Sigma\circ\rr=\Sigma\circ K$ with
$K=\hw\circ\rr$. The exponent in ``$(\Sigma\circ K)^{\pm1}$'' is $+1$, and typing forces it: $\Phi$ and $\Sigma\circ K$
both map $\LR(A/B;c)\to\LR(B^\vee/c;A^\vee)$, while the inverse maps the other way.
\end{corollary}
\begin{proof}
By induction on the slides, Lemma~\ref{lem:B} gives that $X$ and $e_iX$ have identical cell sets after every slide, so
$V(e_iX)=V(X)$. Since $\Sigma$ factors through $V$, $\Sigma(e_iX)=\Sigma(X)$. Now $\hw$ is a composite of $e_i$'s.
\end{proof}

\section{Consequence for H1}
Write $P$ for a mosaic, $T=\kappa(P)\in\LR(\nu/\mu;\lambda)$ (session letters), and $\mathrm{Tao}(P)$ for the readout of
Prop.~5.1(ii). Then
\[
\Phi_G(P)\overset{\text{rot.\ cov.\ (proved)}}{=}R_{iii}(\rho^{-1}P)\overset{\text{G1}}{=}\mathrm{Tao}(\rho^{-1}P)
\overset{\text{G3}}{=}\Phi(\mathrm{Tao}\,P)\overset{\text{Thm~\ref{thm:G4}}}{=}\Sigma(\rr\,\mathrm{Tao}\,P)
\overset{\text{G2 + Lem.~\ref{lem:B}}}{=}\Sigma(\kappa P).
\]
With the factorisation $M=\Phi_G\circ\kappa^{-1}$ (proved last session), this is H1. \textbf{G4 is discharged.} G2 can
be weakened to: $\kappa(P)$ and $\rr(\mathrm{Tao}\,P)$ have the same slide-shape evolution $V$, for example by lying in one
crystal component, which is what K1 says. G3 is needed with direction $\rho^{-1}$. The engine checks the composite
$\Phi(\mathrm{Tao}\,P)=\Phi_G(P)$ directly, 134/134.

\section{Verification}
\begin{center}\begin{tabular}{lrl}\toprule
check & result & script\\\midrule
L0: $\rr(\mathrm{Tao}\,P)$ sits on the cells of $T$ (typing of step~1) & 134/134 & \texttt{g4.py}\\
Lemma~\ref{lem:C} ($\std\rr f=\tilde f$) & 134/134; 13\,171/13\,171 & \texttt{g4.py}, \texttt{g4\_pure.py}\\
$V(Y)=V(\std Y)=V(\hw Y)$ & 134/134 & \texttt{g4.py}\\
Lemma~\ref{lem:B} commutations along $\Sigma$-runs & 606/606 & \texttt{g4.py}\\
Lemma~\ref{lem:D} (outward $=\Sigma$) & 134/134 & \texttt{g4.py}\\
Thm~\ref{thm:G4} and Cor.~\ref{cor:G4}, pure, boxes $\le4\times5$, $|A/B|\le10$, $\mathfrak{gl}_4$ & 13\,171/13\,171 & \texttt{g4\_pure.py}\\
$\Sigma$ constant on whole crystal component & 129\,693 elements / 13\,171 & \texttt{g4\_pure.py}\\
$\Phi$ injective on each type; output LR & 12\,057/12\,057; 13\,171 & \texttt{g4\_pure.py}\\
composite $\Phi(\mathrm{Tao}\,P)=\Phi_G(P)$ (engine) & 134/134 & \texttt{g4.py}\\\midrule
ctrl a: $\Sigma$ of every \emph{other} LR tableau of the type $\ne\Phi(\mathrm{Tao}\,P)$ & fires 134/134; 1003/1003 types & both\\
ctrl b: unrotated $\tilde f$ & fires 134/134 (12 shape-matched) & \texttt{g4.py}\\
ctrl c: tie rule lesioned to ``west'' & fires 23/134 (224 ties met) & \texttt{g4.py}\\\bottomrule
\end{tabular}\end{center}
\textbf{Instrument note.} The first run of \texttt{g4.py} gave L4 at 108/134. The cause was my helper \texttt{inner\_of}, which
dropped rows of $\nu$ lying wholly inside $\kappa$. After the fix all checks read 134/134. The engine's \texttt{switch.sigma} never
had that bug, because it takes the outer shape as an argument.

\section{Gaps}
\begin{enumerate}
\item Lemma~\ref{lem:B} (coplactic) is recalled. It is load-bearing for Cor.~\ref{cor:G4} and for the H1 chain, but
\emph{not} for Thm~\ref{thm:G4}.
\item ``jdt preserves semistandardness'' (used in Lemma~\ref{lem:A}) is recalled. It is a textbook classical fact.
\item Not today's target: G1 (Prop.~5.1), G2 (K1), G3 (Cor.~5.3 with direction $\rho^{-1}$). These are Purbhoo's omitted
proofs and my computed K1. The tie-break interpretation of migration remains a flagged interpretation.
\end{enumerate}

\begin{thebibliography}{9}
\bibitem{P} K.~Purbhoo, \emph{Puzzles, tableaux, and mosaics}, arXiv:0705.1184 (2007).
\end{thebibliography}
\end{document}
